% =====================================================================
%  Tema Beamer Administrativo — șabloane și componente.
%  Se încarcă DUPĂ tema/baza.tex. Folosită de prezentare-comerciala.tex
%  și de prezentare-module.tex: un slide nou = un apel de macro, nu un
%  preambul copiat.
% =====================================================================

% ---------------------------------------------------------------------
% 3. Tema beamer — construită de la zero, fără ornamentele implicite
% ---------------------------------------------------------------------
\usetheme{default}
\setbeamertemplate{navigation symbols}{}
\setbeamercolor{normal text}{fg=brandNavy,bg=white}
\setbeamercolor{structure}{fg=brandBlue}
\setbeamersize{text margin left=0.85cm, text margin right=0.85cm}
\setbeamertemplate{itemize item}{\mbox{\textcolor{brandBlue}{\footnotesize\faIcon{caret-right}}}}
\setbeamertemplate{itemize subitem}{\mbox{\textcolor{brandTeal}{\scriptsize\faIcon{minus}}}}
\setlength{\parskip}{0pt}
% Măsură scurtă peste tot (coloane de 4–8 cm): aliniere la stânga, cu despărțire
% în silabe. Textul justificat lasă găuri urâte pe rânduri atât de scurte.
\AtBeginDocument{\RaggedRight}

% Titlu de slide: titlu gros + linie de accent
\setbeamertemplate{frametitle}{%
  \vspace*{0.34cm}%
  \begin{minipage}{\linewidth}%
    \raggedright
    {\titlufont\bfseries\large\color{brandNavy}\insertframetitle}\par
    \vspace{3pt}
    {\color{brandBlue}\rule{1.7cm}{2.2pt}}%
  \end{minipage}%
  \vspace*{0.14cm}%
}

% Subsol: marcă, poziționare, număr de pagină
\setbeamertemplate{footline}{%
  \hbox{%
    \hspace*{0.85cm}%
    \parbox{\dimexpr\paperwidth-1.7cm\relax}{%
      \textcolor{brandLine}{\rule{\linewidth}{0.5pt}}\par
      \vspace{2pt}
      {\scriptsize\color{brandSlate}{\titlufont ADMINISTRATIVO}\hfill
       ERP administrativ, livrat ca abonament\hfill
       \textbf{\insertframenumber}}%
    }%
  }\vspace*{0.26cm}%
}

% ---------------------------------------------------------------------
% 4. Componente reutilizabile
% ---------------------------------------------------------------------

% Pastilă de status / etichetă colorată
\newcommand{\pastila}[2]{%
  \tikz[baseline=(x.base)]{\node[fill=#1,text=white,rounded corners=2.5pt,
    inner xsep=6pt,inner ysep=2.6pt,font=\bfseries\tiny] (x) {\MakeUppercase{#2}};}}

% Placă de logo — înlocuiește cu \includegraphics{sigla.pdf} când există sigla reală
\newcommand{\logoPlaceholder}[1]{%
  \begin{tikzpicture}[scale=#1,every node/.style={transform shape}]
    \draw[brandTeal,line width=1.1pt,dashed,rounded corners=6pt] (0,0) rectangle (2.6,2.6);
    \node[text=white,font=\titlufont\bfseries\LARGE] at (1.3,1.52) {A};
    \node[text=brandTeal,font=\titlufont\bfseries\tiny] at (1.3,0.62) {LOGO};
  \end{tikzpicture}}

% Card cu iconiță, titlu și text — cărămida întregii prezentări
\newtcolorbox{card}[3][brandBlue]{%
  enhanced, arc=3pt,
  colback=white, colframe=brandLine, boxrule=0.5pt,
  borderline west={2.4pt}{0pt}{#1},
  left=7pt,right=7pt,top=4pt,bottom=4pt,
  halign=flush left, halign title=flush left,
  fonttitle=\bfseries\footnotesize,
  title={\textcolor{#1}{\faIcon{#2}}\hspace{5pt}\textcolor{brandNavy}{#3}},
  colbacktitle=white, coltitle=brandNavy,
  toptitle=1pt, bottomtitle=2pt}

% Cutie de beneficii — fundal deschis, bară de accent
\newtcolorbox{beneficii}[1][Ce câștigă firma]{%
  enhanced, arc=3pt, colback=brandMist, colframe=brandTeal, boxrule=0pt,
  borderline west={3pt}{0pt}{brandTeal},
  left=8pt,right=8pt,top=4pt,bottom=4pt,
  halign=flush left, halign title=flush left,
  fonttitle=\bfseries\footnotesize, coltitle=brandTeal, colbacktitle=brandMist,
  title={\faIcon{check-circle}\hspace{5pt}#1}}

% Cutie de notă comercială evidențiată
\newtcolorbox{nota}[1][Notă]{%
  enhanced, arc=3pt, colback=brandAmber!8, colframe=brandAmber, boxrule=1pt,
  left=9pt,right=9pt,top=6pt,bottom=6pt,
  halign=flush left, halign title=flush left,
  fonttitle=\bfseries\footnotesize, coltitle=white, colbacktitle=brandAmber,
  title={\faIcon{star}\hspace{5pt}#1}}

% Cutie neutră, cu chenar întrerupt — folosită pentru placeholder-uri
\newtcolorbox{placeholder}{%
  enhanced, arc=3pt, colback=brandMist!55, colframe=white, boxrule=0pt,
  borderline={1pt}{0pt}{brandSlate,dashed},
  left=10pt,right=10pt,top=8pt,bottom=8pt}

% Rânduri de listă cu iconiță colorată.
% Fiecare rând ÎȘI ÎNCHIDE singur paragraful (\par la început și la sfârșit):
% altfel `\vspace` de deasupra rămâne în mod orizontal și rândul se lipește de
% textul dinainte. `\hangindent` aliniază continuarea sub text, nu sub iconiță.
%
% Iconița stă într-un `\makebox` de LĂȚIME FIXĂ, egală cu `\hangindent`, și e
% CENTRATĂ în el. Pictogramele Font Awesome au lățimi diferite (`user` are
% 0,875 em, `users` are 1,25 em), deci:
%   - fără casetă fixă, textul pornea din alt punct pe fiecare rând;
%   - cu casetă fixă dar aliniere la stânga, textul se alinia, dar pictogramele
%     rămâneau descentrate una față de alta.
% Alinierea e la STÂNGA, nu centrată. Centrarea pare ideea corectă, dar
% pictogramele Font Awesome au lățimi proprii diferite, deci centrate ies toate
% din altă poziție: marginea din stânga a coloanei devine zimțată, exact ce se
% vedea pe slide-ul Nucleu. Aliniate la stânga, marginile lor formează o linie
% dreaptă, iar diferența de lățime se pierde în spațiul până la text.
%
% Caseta e de 1,25 em (lățimea maximă a unei pictograme Font Awesome), iar
% spațiul până la text e ADĂUGAT SEPARAT, nu lăsat pe seama umpluturii casetei:
% altfel o pictogramă lată umple caseta și textul i se lipește, în timp ce una
% îngustă stă departe. 1,25 + 0,4 = 1,65 em, exact `\hangindent`-ul, deci
% continuarea rândului cade sub prima literă, nu sub pictogramă.
\newcommand{\randlista}[3]{%
  \par\hangindent=1.65em\hangafter=1\noindent
  \makebox[1.25em][l]{\textcolor{#1}{\faIcon{#2}}}\hspace{0.4em}#3\par}
\newcommand{\pct}[2]{\randlista{brandBlue}{#1}{#2}}
\newcommand{\pctv}[2]{\randlista{brandGreen}{#1}{#2}}

% Titlu de secțiune peste ecran plin: iconiță, titlu, subtitlu
\newcommand{\separator}[3]{%
\begin{frame}[plain]
  \begin{tikzpicture}[remember picture,overlay]
    \fill[brandNavy] (current page.south west) rectangle (current page.north east);
    \fill[brandTeal,opacity=0.10] ([xshift=-3cm,yshift=-2.4cm]current page.east) circle (5.2cm);
    \fill[brandBlue,opacity=0.16] ([xshift=-1.2cm,yshift=2.2cm]current page.east) circle (3.4cm);
    \node[anchor=west,align=left] at ([xshift=1.6cm]current page.west) {%
      \begin{minipage}{10.5cm}
        \raggedright
        {\color{brandTeal}\Large\faIcon{#1}}\par\vspace{9pt}
        {\titlufont\bfseries\huge\color{white} #2}\par\vspace{8pt}
        {\color{brandLine}\small #3}
      \end{minipage}};
  \end{tikzpicture}
\end{frame}}


% ---------------------------------------------------------------------
% 5. Tiparul unui slide de modul
% ---------------------------------------------------------------------
% Toată secțiunea de module e construită dintr-un singur macro, ca cele
% douăzeci de slide-uri să nu ajungă douăzeci de layout-uri ușor diferite.
% Un modul nou = un apel nou, nu un slide copiat și ajustat.
%
%   #1 titlul slide-ului
%   #2 grupul din meniu, tipărit ca a doua pastilă (Personal, Financiar, ...)
%   #3 paragraful de deschidere — ce e modulul, în două-trei rânduri
%   #4 funcționalitățile, ca rânduri \pct{iconiță}{text}
%   #5 beneficiile, ca rânduri \pctv{check}{text}
%
% `\vspace{-2pt}` de la început nu e cosmetic: corpul cadrului începe cu un
% grup `{...}`, iar beamer l-ar înghiți ca SUBTITLU (vezi README, capcana 1).
\newcommand{\modul}[5]{%
\begin{frame}{#1}
  \vspace{-2pt}
  {\pastila{brandGreen}{disponibil}\hspace{5pt}\pastila{brandBlue}{#2}}
  \vspace{7pt}
  \begin{columns}[T,onlytextwidth]
    \begin{column}{0.52\textwidth}
      {\footnotesize #3}
      \vspace{5pt}
      {\footnotesize #4}
    \end{column}
    \begin{column}{0.44\textwidth}
      \begin{beneficii}[Ce câștigă firma]
        \footnotesize #5
      \end{beneficii}
    \end{column}
  \end{columns}
\end{frame}}

% Card compact — doar pentru pagina cu toate modulele, unde încap opt rânduri într-o
% coloană. Aceleași culori ca `card`, dar cu jumătate din padding.
\newtcolorbox{cardm}[3][brandBlue]{%
  enhanced, arc=3pt,
  colback=white, colframe=brandLine, boxrule=0.5pt,
  borderline west={2.4pt}{0pt}{#1},
  left=6pt,right=6pt,top=1pt,bottom=2pt,
  halign=flush left, halign title=flush left,
  fonttitle=\bfseries\scriptsize,
  title={\textcolor{#1}{\faIcon{#2}}\hspace{4pt}\textcolor{brandNavy}{#3}},
  colbacktitle=white, coltitle=brandNavy,
  toptitle=0pt, bottomtitle=1pt}

% Varianta pentru un modul care ÎNCĂ NU e livrat. Singura diferență față de
% `\modul` e pastila: chihlimbar „în dezvoltare" în locul verdelui „disponibil".
% Deliberat un macro separat, nu un argument opțional: culoarea pastilei e
% singurul lucru care spune clientului că plătește pentru ceva ce încă nu poate
% folosi azi, deci nu trebuie să depindă de un parametru ușor de uitat.
\newcommand{\modulviitor}[5]{%
\begin{frame}{#1}
  \vspace{-2pt}
  {\pastila{brandAmber}{în dezvoltare}\hspace{5pt}\pastila{brandBlue}{#2}}
  \vspace{7pt}
  \begin{columns}[T,onlytextwidth]
    \begin{column}{0.52\textwidth}
      {\footnotesize #3}
      \vspace{5pt}
      {\footnotesize #4}
    \end{column}
    \begin{column}{0.44\textwidth}
      \begin{beneficii}[Ce va câștiga firma]
        \footnotesize #5
      \end{beneficii}
    \end{column}
  \end{columns}
\end{frame}}

% Celulă din pagina cu toate modulele: iconiță + denumire, pe un rând.
%
% `\tiny` se deschide ÎNAINTE de paragraf și `\par` se închide ÎN INTERIORUL
% grupului, deliberat. TeX fixează `\baselineskip` la `\par`, din mărimea
% curentă: cu `{\tiny #2}` doar pe text, literele erau mici, dar rândurile
% rămâneau la interlinia de 11pt — opt module ocupau cât un slide întreg.
\newcommand{\celula}[2]{%
  {\tiny\par\hangindent=1.6em\hangafter=1\noindent
   \makebox[1.25em][l]{\textcolor{brandBlue}{\faIcon{#1}}}\hspace{0.35em}#2\par}}

% Rând de modul cu preț: denumirea la stânga, prețul lipit de marginea dreaptă.
% `\makebox` de lățimea coloanei aduce modul LR, singurul în care `\hfill` chiar
% împinge — într-un paragraf cu `\RaggedRight` activ, glue-ul din dreapta l-ar
% fi înghițit, iar prețurile ar fi stat lipite de denumire.
\newcommand{\modulpret}[2]{%
  {\tiny\par\noindent\makebox[\linewidth][l]{#1\hfill\textbf{#2}}\par}}

% Cifră mare cu etichetă dedesubt — banda de sinteză de pe slide-ul 1.
\newcommand{\cifra}[2]{%
  \begin{minipage}[t]{3.15cm}
    \centering
    {\titlufont\bfseries\Large\color{brandBlue}#1}\par\vspace{1pt}
    {\scriptsize\color{brandSlate}#2}
  \end{minipage}}

% ---------------------------------------------------------------------
% 6. Slide de modul cu captură din aplicație
% ---------------------------------------------------------------------
% Fereastra în care stă captura: bară de titlu cu trei puncte, ca un ecran
% văzut „prin geam" — același motiv ca pe paginile de modul ale site-ului.
% Fără margine interioară: chenarul atinge imaginea, altfel captura pare
% lipită pe o foaie, nu deschisă pe un ecran.
\newtcolorbox{fereastra}{%
  enhanced, arc=3pt, colback=white, colframe=brandLine, boxrule=0.5pt,
  left=0pt,right=0pt,top=0pt,bottom=0pt, boxsep=0pt,
  halign=center,
  colbacktitle=brandMist, coltitle=brandLine,
  toptitle=1.2pt, bottomtitle=1.2pt, lefttitle=4pt,
  fonttitle=\fontsize{3.5}{4}\selectfont,
  title={\faIcon{circle}\,\faIcon{circle}\,\faIcon{circle}}}

% Banda de beneficii de sub coloane: aceleași culori ca `beneficii`, dar
% culcată pe toată lățimea, cu beneficiile pe trei coloane.
%
% De ce nu cutia `beneficii` sub captură, în coloana dreaptă: prima variantă
% așa era, și ieșea peste subsol cu 25–50 pt pe fiecare slide. Captura
% 16:10 ocupă aproape toată înălțimea coloanei; beneficiile scurte, puse în
% lat, costă un singur rând în loc de patru.
\newtcolorbox{bandabeneficii}{%
  enhanced, arc=3pt, colback=brandMist, colframe=brandTeal, boxrule=0pt,
  borderline west={3pt}{0pt}{brandTeal},
  left=8pt,right=8pt,top=3pt,bottom=4pt}

% Un beneficiu din bandă: o treime din lățime, aliniat sus. `\hfill` de la
% final împarte golul rămas între cele trei coloane.
%
% `\RaggedRight` se repune ÎN minipage: minipage reface alinierea implicită
% (`\@parboxrestore`), deci textul ieșea justificat pe o coloană de 4 cm —
% câte trei `Underfull \hbox` pe fiecare slide.
\newcommand{\benef}[1]{%
  \begin{minipage}[t]{0.27\linewidth}\RaggedRight\pctv{check}{#1}\end{minipage}\hfill}

% Același tipar ca `\modul`, cu o coloană dreaptă diferită: captura, în
% fereastră, iar beneficiile într-o bandă pe toată lățimea, dedesubt.
%
%   #1–#4 exact ca la `\modul`
%   #5    EXACT trei beneficii, ca `\benef{...}` — nu `\pctv`, și fără
%         `\par` între ele: un `\par` le-ar așeza unul sub altul
%   #6    conținutul ferestrei: de regulă `\captura{cheie}`; Portalul pune
%         aici două capturi de telefon alăturate
%
% Pastila de status NU e parametru, din același motiv ca la `\modulviitor`.
%
% Eticheta „Ce câștigă firma" stă în prima coloană a benzii, nu deasupra ei:
% pe un rând separat costa ~10 pt, exact cât ieșea banda peste subsol pe
% slide-urile cu șase funcționalități.
\newcommand{\modulcaptura}[6]{%
\begin{frame}{#1}
  \vspace{-2pt}
  {\pastila{brandGreen}{disponibil}\hspace{5pt}\pastila{brandBlue}{#2}}
  \vspace{5pt}
  \begin{columns}[T,onlytextwidth]
    \begin{column}{0.5\textwidth}
      {\footnotesize #3}
      \vspace{4pt}
      {\footnotesize #4}
    \end{column}
    \begin{column}{0.47\textwidth}
      \begin{fereastra}#6\end{fereastra}
    \end{column}
  \end{columns}
  \vspace{4pt}
  \begin{bandabeneficii}
    \scriptsize
    \begin{minipage}[t]{0.14\linewidth}\RaggedRight
      \randlista{brandTeal}{check-circle}{\bfseries\color{brandTeal}Ce câștigă firma}
    \end{minipage}\hfill
    #5
  \end{bandabeneficii}
\end{frame}}

% Captura de birou a unui modul, după cheia din `src/config/features.ts`.
% Înălțimea e plafonată, nu lățimea: capturile sunt 16:10, iar la lățimea
% coloanei ar împinge banda de beneficii peste subsol.
\newcommand{\captura}[1]{%
  \includegraphics[width=\linewidth,height=3.0cm,keepaspectratio]{capturi/#1}}
